\section{The non-induced imaginary-quadratic physical-projection obstruction}
\label{sec:imaginary-projection-obstruction}

\paragraph{Status and scope.}
The complete author reading and two independent whole-proof
readings have passed; their actual check scopes are recorded separately.  It uses the actual non-induced cubic character
already fixed in MA-CH-003.  The different prime is retained in $S$,
and the final input belongs to the original coordinate-flat source.
Only direct membership in the full source PW range is needed.
The modified finite test vector below is not identified with the
primitive canonical vector or with a new Artin character.

Keep the established arithmetic data
\[
 \begin{gathered}
 K=\mathbb Q(\zeta_3),\quad E=K(269^{1/3}),\quad
 \psi=\psi_{269,3},\quad
 \mathfrak p=(\sqrt{-3}),\quad \mathfrak q=(269),\\
 N\mathfrak p=3,\quad q=N\mathfrak q=269^2,\quad
 \mathfrak f_\psi=\mathfrak q,\quad Q=3q,\quad
 S=\{\infty,\mathfrak p,\mathfrak q\},\qquad S'=S.
 \end{gathered}
 \tag{MA-PI.120}
\]
Here $Q$ is the primitive analytic conductor.
The class number is one, $|\mathcal O_K^\times|=6$,
and $\psi_{\mathfrak p}$ is trivial: the prime over $3$
splits in $E/K$.  The local different exponent at
$\mathfrak p$ is one.  At $\mathfrak q$ the character
has conductor exponent one and no inertia invariants.
These facts, including $\psi^c=\psi^{-1}\ne\psi$ and
non-descent from a Dirichlet norm character, are the
already proved MA-CH-003/MA-PF-003/005 facts.

Use the coherent additive characters of MA-CGP-001.
At infinity the Fourier kernel is
$e^{-2\pi i\operatorname{Tr}_{\mathbb C/\mathbb R}(z\xi)}$,
with self-dual measure $2\,d(\operatorname{Re}z)\,d(\operatorname{Im}z)$.
The normalized complex absolute value is $|z|_{\mathbb C}=|z|^2$.
Put $\pi=\sqrt{-3}$, and take the finite functions
\[
 \begin{split}
 h_{\mathfrak p}
   &=\mathbf1_{\mathcal O_{\mathfrak p}^{\times}}
            -3\mathbf1_{\pi\mathcal O_{\mathfrak p}^{\times}}
     =\mathbf1_{\mathcal O_{\mathfrak p}}
            -4\mathbf1_{\pi\mathcal O_{\mathfrak p}}
            +3\mathbf1_{\pi^2\mathcal O_{\mathfrak p}},\\
 \kappa_{\mathfrak q}
   &=\psi_{\mathfrak q}^{-1}|_{\mathcal O_{\mathfrak q}^{\times}}\,
                             \mathbf1_{\mathcal O_{\mathfrak q}^{\times}},\\
 \varphi_f&=f\otimes h_{\mathfrak p}\otimes\kappa_{\mathfrak q}
                    \quad\text{on }\mathbb A_{K,S}.
 \end{split}
 \tag{MA-PI.121}
\]
The global extension has $\mathbf1_{\mathcal O_v}$ outside $S$.
Let $\mathcal T_S$ denote its actual centered, fixed-$\psi$
global theta map, and put
$Y_f=\mathcal T_S(\varphi_f)$ for radial $f\in\mathcal S(\mathbb C)$.
The finite ramified factor has value and integral zero,
so all these global data satisfy both polar conditions,
without conditions on $f(0)$ or its integral.

For $Y>0$ let $\rho_Y$ be a real smooth profile with
\[
 0\le\rho_Y\le1,\qquad
 \rho_Y(r)=1\ (r\le Y/2),\qquad
 \rho_Y(r)=0\ (r\ge2Y),\qquad
 \mathfrak B_{S,Y}
      =\mathcal F_S^{-1}M_{\rho_Y(|\xi|_S)}\mathcal F_S.
 \tag{MA-PI.122}
\]
For either of the independently specified metrics
\[
 \begin{aligned}
 H^{\rm sym}_{a,s}
    &=L^2(\mathbb R,e^{2a|x|}(1+x^2)^{s/2}dx),
                 &&0<a<1/2,\ s\ge0,\\
 H^{\rm one}_{b,s}
    &=L^2(\mathbb R,e^{2bx}(1+x^2)^{s/2}dx),
                 &&|b|<1/2,\ s\ge0,
 \end{aligned}
 \tag{MA-PI.123}
\]
write $Q^H_{S,Y}$ for the orthogonal projection onto the
closure of the full actual invariant global PW image,
with the source product Fourier bound $Y$ and $S'=S$.
This is the corrected centered realization in the first
metric.  Its exact printed realization in the second
metric is identified below.  The projection is onto the
full original source level, not only the image of the
particular family (MA-PI.121).

\paragraph{MA-PI-004 (a non-induced character witness in the original domains).}
For either fixed metric in (MA-PI.123), there is $Y_0>0$
such that, for every $Y\ge Y_0$ and every profile
(MA-PI.122), there is a radial Schwartz function $f_Y$
identically zero near the complex origin for which
\[
 \varphi_{f_Y}\in V_S,\qquad
 Q^H_{S,Y}Y_{f_Y}
       \ne\mathcal T_S(\mathfrak B_{S,Y}\varphi_{f_Y}).
 \tag{MA-PI.124}
\]
Here $V_S$ is the original inherited coordinate-flat LF
source of MA-GA.2.  The cutoff output is actual
$S$-local Schwartz data satisfying both polar conditions.
Thus $Y_{f_Y}$ and
$\mathcal T_S(\mathfrak B_{S,Y}\varphi_{f_Y})$
belong to $\mathcal E$.
The set $S$ includes the different and the conductor and
satisfies the arithmetic hypotheses of the original
source setup, including Theorem~12.4.1.1, p.~514.
No Hecke zero-location premise is used.

\begin{proof}
\emph{1. The exact finite Fourier shells and unchanged physical cutoff.}
The self-dual volume at $\mathfrak p$ is
$v=\operatorname{vol}(\mathcal O_{\mathfrak p})=3^{-1/2}$.
Annihilators of balls, with different exponent one, give
\[
 \widehat h_{\mathfrak p}(\xi)
   =v\left(\mathbf1_{\{|\xi|_{\mathfrak p}\le3\}}
       -\frac43\mathbf1_{\{|\xi|_{\mathfrak p}\le9\}}
       +\frac13\mathbf1_{\{|\xi|_{\mathfrak p}\le27\}}\right)
   =3^{-3/2}h_{\mathfrak p}(\pi^3\xi).
 \tag{MA-PI.125}
\]
Indeed $\widehat{\mathbf1_{\pi^j\mathcal O}}=
3^{-j}v\,\mathbf1_{\pi^{-j-1}\mathcal O}$, for $j=0,1,2$.
The values in (MA-PI.125) are $-v$ on the shell of norm
$9$, $v/3$ on the shell of norm $27$, and zero elsewhere.
The original $h_{\mathfrak p}$ is supported away from zero
and has integral $v-4v/3+3v/9=0$.

The primitive finite Gauss calculation of MA-CQ-001,
with conductor exponent one and different exponent zero,
gives a single Fourier shell
$|\xi_{\mathfrak q}|_{\mathfrak q}=q$ for
$\kappa_{\mathfrak q}$, with a nonzero Gauss factor of
modulus $q^{-1/2}$.  Its exact phase cancels when the
actual inverse Fourier transform is applied.
The finite product Fourier norms of (MA-PI.121) are
therefore exactly $9q$ and $27q$.

For later use split (MA-PI.125) into its two shells.
Their actual inverse Fourier transforms are
\[
 h_9=\mathbf1_{\mathcal O_{\mathfrak p}}
                  -3\mathbf1_{\pi\mathcal O_{\mathfrak p}},
 \qquad
 h_{27}=3\mathbf1_{\pi^2\mathcal O_{\mathfrak p}}
                  -\mathbf1_{\pi\mathcal O_{\mathfrak p}},
 \qquad h_9+h_{27}=h_{\mathfrak p}.
 \tag{MA-PI.126}
\]
For example
$\mathcal F_{\mathfrak p}^{-1}\mathbf1_{\{|\xi|\le3^k\}}
 =3^kv\,\mathbf1_{\pi^{k-1}\mathcal O_{\mathfrak p}}$.
Multiplying the shell differences by $-v$ and $v/3$
proves (MA-PI.126), including every constant.
Consequently the physical cutoff has the literal decomposition
\[
 \begin{split}
 \mathfrak B_{S,Y}\varphi_f
   &=\sum_{r\in\{9,27\}}
          f_{Y,r}\otimes h_r\otimes\kappa_{\mathfrak q},\\
 f_{Y,r}
   &=\mathcal F_{\mathbb C}^{-1}
           \bigl(\rho_Y(rq|z|_{\mathbb C})\widehat f(z)\bigr).
 \end{split}
 \tag{MA-PI.127}
\]
This computation retains the two finite Fourier bands;
it does not replace them by a single cutoff of a different
level.  Every summand is Schwartz and has both global
polar moments zero because of $\kappa_{\mathfrak q}$.
Conjugating coherent Fourier signs only reflects norm
coordinates, leaving these norm cutoffs unchanged.

\emph{2. Actual theta, its finite filter, and its effective reflection scale.}
Let $Y_f^{\rm prim}$ denote the primitive-vector theta
of MA-CGP-001.  With $f$ radial its exact formula is
\[
 Y_f^{\rm prim}(x)
    =6e^{x/2}\sum_{\mathfrak a\ne0}
             \psi(\mathfrak a)
                       f\bigl((N\mathfrak a\,e^x)^{1/2}\bigr).
 \tag{MA-PI.128}
\]
The ideal coefficient is extended by zero at $\mathfrak q$.
Class number one and the six units give the factor $6$.
The equality $\psi_{\mathfrak p}=1$ makes the actual
finite change in (MA-PI.121) exactly
\[
 Y_f(x)=Y_f^{\rm prim}(x)
          -\frac4{\sqrt3}Y_f^{\rm prim}(x+\log3)
          +Y_f^{\rm prim}(x+2\log3).
 \tag{MA-PI.129}
\]
There is no unrecorded uniformizer phase.  The principal-idele
relation on $\pi$ forces $\psi_{\mathfrak q}(\pi)=1$:
the archimedean component is trivial, the component at
$\mathfrak p$ is trivial, and all other finite components
are unramified on the unit $\pi$.  Any complex rotation
from multiplication by $\pi$ is invisible to radial $f$.
The factor $3^{-j/2}$ from the centered shift accounts for
the two coefficients in (MA-PI.129).

The actual Poisson formula for the primitive finite vector,
or (MA-PI.125) together with its ramified Gauss factor,
therefore yields
\[
 Y_f(x)=\epsilon_\psi\,Y^{\,\bar\psi}_{\widehat f}
                     (-x-\log Q_*),\qquad
 Q_*=27q=9Q,\qquad |\epsilon_\psi|=1.
 \tag{MA-PI.130}
\]
Here the superscript means that the same $h_{\mathfrak p}$
is used and the ramified character is conjugated.
For a direct check using (MA-PI.129), the primitive
reflection is at $-\log Q$.  Replacing $x$ by
$-x-\log(9Q)$ in the three dual shifted terms reverses
their order; their coefficients $1,-4/\sqrt3,1$ are
unchanged.  This proves (MA-PI.130), with the same phase.
Thus $Q_*$ is an effective scale of this test vector,
not a replacement for the primitive conductor $Q=3q$.

The direct ideal sum at $x\to+\infty$ and the reflected
sum at $x\to-\infty$ prove
$Y_f\in\mathcal E$, continuously for radial Schwartz $f$.
For completeness, a Schwartz bound of order $2M$ on the
archimedean variable bounds the direct sum by a constant
times $e^{-(M-1/2)x}\sum_{\mathfrak a}(N\mathfrak a)^{-M}$,
for $M>1$.  Every norm derivative satisfies the same kind
of estimate.  The reflected identity and continuity of
Fourier transform on Schwartz space give the other end.
This proves continuity by finitely many source seminorms
for every prescribed target seminorm.  It applies uniformly
to bounded Schwartz families.

Write $F(t)=f(\sqrt t)$ for the radial norm profile, and
$M_F(s)=\int_0^\infty F(t)t^{s-1}dt$.
Actual unfolding in $\Re s>1$ gives
\[
 \begin{split}
 \mathcal M Y_f(s-\tfrac12)&=6A_\psi(s)M_F(s),\\
 A_\psi(s)&=E_3(s)L_K(s,\psi),\\
 E_3(s)&=(1-3^{-s})(1-3^{1-s})
           =1-4\,3^{-s}+3\,3^{-2s}.
 \end{split}
 \tag{MA-PI.131}
\]
The factor replaces the primitive local geometric series
at $\mathfrak p$ by the weights $1,-3,0,\ldots$.
One has exactly
$E_3(s)=3^{1-2s}E_3(1-s)$.

The Gaussian $f_0(z)=e^{-2\pi|z|^2}$ is self-dual for the
declared complex Fourier measure, and
$M_{F_0}(s)=(2\pi)^{-s}\Gamma(s)=\Gamma_{\mathbb C}(s)/2$,
where $\Gamma_{\mathbb C}(s)=2(2\pi)^{-s}\Gamma(s)$.
Taking its Mellin transform in (MA-PI.130) proves
\[
 \Gamma_{\mathbb C}(s)A_\psi(s)
   =\epsilon_\psi Q_*^{1/2-s}
                \Gamma_{\mathbb C}(1-s)A_{\bar\psi}(1-s).
 \tag{MA-PI.132}
\]
Both completed expressions are entire, by the
all-exponential theta bounds.  Multiplication by the
entire reciprocal gamma function also proves that
$A_\psi$ and $A_{\bar\psi}$ are entire.  The effective
scale in (MA-PI.132) is derived from the actual source,
not assigned as an arithmetic conductor.
For arbitrary radial Schwartz $f$, $M_F(s)$ is holomorphic
in $\Re s>0$.  Analytic continuation therefore extends
(MA-PI.131) to that half-plane, which includes every
line used in the Gram calculation.

\emph{3. The exact radial Gram form and its temperateness.}
Fix $b\in(-1/2,1/2)$ and put
$\sigma=1/2+b$, $d=1-\sigma$.
Given $A_u\in C_c^\infty(\mathbb R)$, define a radial
frequency vector by
\[
 u(z)=U(|z|_{\mathbb C}),\qquad
 U(t)=t^{\sigma-1}A_u(\log t),\qquad
 f=\mathcal F_{\mathbb C}^{-1}u.
 \tag{MA-PI.133}
\]
It is smooth with compact support away from the complex
origin; $f$ is radial Schwartz.
The radial local Fourier equation is
\[
 M_F(s)=\frac{\Gamma_{\mathbb C}(s)}
                  {\Gamma_{\mathbb C}(1-s)}M_U(1-s),
                 \qquad0<\Re s<1.
 \tag{MA-PI.134}
\]
This is Tate's local equation, Theorem~2.4.1, with the
angular index zero in the complex table, thesis p.~(2.16),
PDF p.~23.  The table gives
$(2\pi)^{1-2s}\Gamma(s)/\Gamma(1-s)$.
Its angular multiplicative-measure factor $2\pi$ occurs
on both sides and cancels for our norm-profile Mellin
integrals.  The self-dual Gaussian above fixes the phase
as one.  On radial functions the inverse and forward
Fourier transforms agree, since their square is reflection.
These observations verify (MA-PI.134) in the declared
normalization rather than importing an unspecified gamma factor.

For $\widehat A_u(t)=\int A_u(r)e^{-itr}dr$ one has
$M_U(1-\sigma-it)=\widehat A_u(t)$.
Equations (MA-PI.131)--(MA-PI.134) and Mellin Plancherel
therefore give, for $f,h$ associated to $u,v$,
\[
 \begin{split}
 \int_{\mathbb R}Y_f(x)\overline{Y_h(x)}e^{2bx}dx
   &=\frac1{2\pi}\int_{\mathbb R}
         W_b(t)\widehat A_u(t)\overline{\widehat A_v(t)}dt,\\
 W_b(t)&=36Q_*^{-2b}
             A_\psi(d+it)A_{\bar\psi}(d-it).
 \end{split}
 \tag{MA-PI.135}
\]
For real $t$ this symbol is real and nonnegative, because
$A_{\bar\psi}(\bar s)=\overline{A_\psi(s)}$.
That identity follows first from the absolutely convergent
ideal series and the real-coefficient $E_3$, and then by
analytic continuation.

Here is an actual-source proof that $W_b(t)\,dt$ is tempered;
no additional vertical-growth theorem is needed.
Fix a nonnegative nonzero $A_0\in C_c^\infty(\mathbb R)$.
Its Fourier transform has modulus at least $c>0$ on
$[-\delta,\delta]$ for some $\delta>0$.
For each $T\in\mathbb R$ use
$A_T(r)=e^{iTr}A_0(r)$ in (MA-PI.133).
The resulting $u_T$ have support in one fixed compact
annulus.  The chain rule shows that each fixed Schwartz
seminorm of $u_T$ is at most $C(1+|T|)^N$: powers and
derivatives of $|z|_{\mathbb C}$ and its logarithm are
bounded on that annulus, and each derivative of the
modulation contributes at most a power of $T$.
Fourier continuity gives the same assertion for any
fixed seminorm of $f_T$.  The actual theta continuity
proved after (MA-PI.130), followed by the continuous
inclusion $\mathcal E\to L^2(e^{2bx}dx)$, consequently gives
$\|Y_{f_T}\|^2\le C'(1+|T|)^{N'}$ for fixed $C',N'$.
Applying the positive identity (MA-PI.135) and
$\widehat A_T(t)=\widehat A_0(t-T)$ yields
\[
       \int_{T-\delta}^{T+\delta}W_b(t)dt
                         \le C''(1+|T|)^{N'}.
 \tag{MA-PI.136}
\]
A covering by intervals of this fixed length now gives
$\int W_b(t)(1+|t|)^{-M}dt<\infty$ for sufficiently
large $M$.  Thus it is a positive tempered measure and
its inverse Fourier transform $\kappa_b$ is defined
as a tempered distribution.

\emph{4. Noncompact support, with Hermitian rather than assumed even symmetry.}
The symbol has the entire continuation
\[
       \mathscr W_b(z)
         =36Q_*^{-2b}A_\psi(d+iz)A_{\bar\psi}(d-iz).
 \tag{MA-PI.137}
\]
It cannot satisfy an exponential-type bound.  At
$s_n=-n-1/2$, (MA-PI.132) and the gamma reflection
identity give
\[
 |A_\psi(s_n)|
   =Q_*^{n+1}(2\pi)^{-2n-2}
       \frac{\Gamma(n+3/2)^2}{\pi}
                            |A_{\bar\psi}(n+3/2)|.
 \tag{MA-PI.138}
\]
Indeed
$\Gamma(-n-1/2)\Gamma(n+3/2)=(-1)^{n+1}\pi$.
Absolute convergence of the ideal series at $2$ and
dominated convergence show that $L_K(t,\bar\psi)\to1$
as the real variable $t\to+\infty$.  Also $E_3(t)\to1$.
Both required positive-axis $A$ factors therefore have
modulus at least $1/2$ for all sufficiently large $n$.
With $T_n=n+1/2+d$, (MA-PI.137)--(MA-PI.138) imply
\[
 |\mathscr W_b(iT_n)|
   \ge c_b\,Q_*^{n+1}(2\pi)^{-2n-2}\Gamma(n+3/2)^2.
 \tag{MA-PI.139}
\]
Since $\Gamma(n+3/2)\ge(\sqrt\pi/2)n!$, this grows
faster than every polynomial times $e^{R T_n}$.
If $\kappa_b$ had compact support, the finite-order
distribution estimate applied to a cutoff times
$e^{-izx}$ would give precisely such a bound for its
entire Fourier transform, as proved in (MA-PI.18).
Agreement on the real axis identifies that transform
with (MA-PI.137), a contradiction.

Although no evenness of $W_b$ is assumed, its reality
on the real axis gives the Hermitian distribution identity
$\kappa_b(-x)=\overline{\kappa_b(x)}$.
In particular its support is symmetric.  The separated
test argument from MA-PI-001 therefore applies unchanged:
if every compact smooth high/low log-frequency pair
beyond a fixed gap had zero Gram pairing, approximate
a delta in the high variable.  This would make
$\kappa_b$ vanish on the corresponding negative
half-line; Hermitian symmetry would give vanishing
on the positive half-line too.  That would make its
support compact.  A nonzero limiting distribution
pairing has a nonzero ordinary smooth approximate-delta
member, so actual compact test functions suffice.

Choose the gap to obtain radial Schwartz $f,h$ with
compact smooth Fourier supports and
\[
 \begin{gathered}
 \operatorname{supp}\widehat f
       \subset\{|z|_{\mathbb C}>2/(9q)\},\qquad
 \operatorname{supp}\widehat h
       \subset\{0<|z|_{\mathbb C}<1/(54q)\},\\
 D_b:=\int Y_f(x)\overline{Y_h(x)}e^{2bx}dx\ne0.
 \end{gathered}
 \tag{MA-PI.140}
\]
No assertion is made for every preassigned pair of
bounded annuli.

\emph{5. Direct PW membership and actual dilation in both metrics.}
For $L>0$ put $Y=e^L$ and use the complex unitary
frequency dilation
\[
 u_L(z)=e^{-L/2}u(e^{-L/2}z),\qquad
 f_L(z)=e^{L/2}f(e^{L/2}z),
 \tag{MA-PI.141}
\]
and similarly for $v,h$.  Complex norm frequencies
are multiplied by $e^L$.  Exact theta summation gives
$Y_{f_L}(x)=Y_f(x+L)$ and the same formula for $h$.
The smallest finite Fourier norm is $9q$ and the
largest is $27q$.  Equations (MA-PI.122) and
(MA-PI.140) therefore imply
\[
 \mathfrak B_{S,Y}\varphi_{f_L}=0,\qquad
 \mathfrak B_{S,Y}\varphi_{h_L}=\varphi_{h_L},\qquad
 Y_{h_L}\in\operatorname{Ran}Q^H_{S,Y}.
 \tag{MA-PI.142}
\]
For the last inclusion the entire $S$-Fourier support
is compact with product norm $<Y/2<Y$.
Both polar conditions hold at $\mathfrak q$.
This is direct membership in Masullo's original
Definition~7.3.2.3 with $S'=S$, followed by its
actual global image and Hilbert closure.  It does not
require the modified family to exhaust that range.

For the symmetric metric choose $b=-a$ when constructing
(MA-PI.140); for the one-sided metric use its own $b$.
The exact two changes of variable give respectively
\[
 \begin{split}
 e^{-2aL}L^{-s}
     \langle Y_{f_L},Y_{h_L}\rangle_{H^{\rm sym}_{a,s}}
                    &\longrightarrow D_{-a}\ne0,\\
 e^{2bL}L^{-s}
     \langle Y_{f_L},Y_{h_L}\rangle_{H^{\rm one}_{b,s}}
                    &\longrightarrow D_b\ne0.
 \end{split}
 \tag{MA-PI.143}
\]
The polynomial ratio is
$(1+(x-L)^2)^{s/2}/L^s$, bounded for $L\ge1$ by
$C_s(2+|x|)^s$ and tending to one.
In the first line the normalized exponential factor
tends to $e^{-2ax}$ and is bounded by $e^{2a|x|}$;
in the second it is exactly $e^{2bx}$.
The actual $\mathcal E$ bounds prove domination.
Thus the pairings are nonzero for all sufficiently
large $Y$, uniformly over the permitted profiles.
Orthogonality and (MA-PI.142) give a nonzero pairing
of $Q^H_{S,Y}Y_{f_L}
-\mathcal T_S(\mathfrak B_{S,Y}\varphi_{f_L})$
against the actual range vector $Y_{h_L}$.

\emph{6. A finite flat witness, including the two-shell cutoff output.}
Fix this finite $Y$ and profile.  Choose a radial
$\omega\in C_c^\infty(\mathbb C)$ equal to one near zero,
and put
\[
 f_{L,\varepsilon}(z)
       =(1-\omega(z/\varepsilon))f_L(z),\qquad
 r_\varepsilon(z)=\omega(z/\varepsilon)f_L(z)
                        =g_\varepsilon(z/\varepsilon).
 \tag{MA-PI.144}
\]
Here $g_\varepsilon(w)=\omega(w)f_L(\varepsilon w)$
is a bounded family of radial compact smooth functions.
No value or derivative vanishing of $f_L$ is needed:
the finite ramified factor already removes both polar
terms.  The original input
$\varphi_{f_{L,\varepsilon}}$ is coordinate-flat at
every place of $S$.

The actual theta series gives exactly
$Y_{r_\varepsilon}(x)=
\varepsilon Y_{g_\varepsilon}(x-2\log\varepsilon)$.
Uniform $\mathcal E$ bounds then yield
\[
 \begin{split}
 \|Y_{r_\varepsilon}\|_{H^{\rm sym}_{a,s}}
    &\le C\varepsilon^{1-2a}
                    (1+|\log\varepsilon|)^{s/2}\longrightarrow0,\\
 \|Y_{r_\varepsilon}\|_{H^{\rm one}_{b,s}}
    &\le C\varepsilon^{1+2b}
                    (1+|\log\varepsilon|)^{s/2}\longrightarrow0.
 \end{split}
 \tag{MA-PI.145}
\]
For the second line substitute $x=y+2\log\varepsilon$
in the squared norm; it gives the exact factor
$\varepsilon^{2+4b}$.  For the first line use
$e^{2a|y+2\log\varepsilon|}
 \le\varepsilon^{-4a}e^{2a|y|}$.
All asserted exponents are positive.

For every real-coordinate multi-index $\alpha$,
differentiation of the complex Fourier integral gives
\[
 \sup_\xi|\partial_\xi^\alpha\widehat r_\varepsilon(\xi)|
                         \le C_\alpha\varepsilon^{|\alpha|+2}.
 \tag{MA-PI.146}
\]
The power $2$ is the real dimension of $\mathbb C$;
the self-dual factor $2$ and Fourier constants only
change $C_\alpha$.  For each of the two fixed shells,
multiplication by $\rho_Y(rq|z|_{\mathbb C})$ gives
compact smooth Fourier functions tending to zero
in every smooth seminorm.  Their inverse transforms
therefore tend to zero in radial Schwartz space.

Apply the exact decomposition (MA-PI.127) to this error.
For either fixed finite factor $h_9$ or $h_{27}$ its
actual theta map is continuous from radial Schwartz
space to $\mathcal E$.  One can see this directly
from (MA-PI.126): its image is respectively
$Y_f^{\rm prim}-\sqrt3Y_f^{\rm prim}(\,\cdot+\log3)$
or
$Y_f^{\rm prim}(\,\cdot+2\log3)
 -3^{-1/2}Y_f^{\rm prim}(\,\cdot+\log3)$.
The ramified factor remains unchanged and removes
both polar moments in each summand.  Consequently
\[
 \|\mathcal T_S(\mathfrak B_{S,Y}\varphi_{r_\varepsilon})\|_H
                                      \longrightarrow0.
 \tag{MA-PI.147}
\]
The profile and its derivative constants were fixed
before choosing $\varepsilon$; no uniform family
derivative bound has been assumed.

Use boundedness of the actual orthogonal projector and
test the difference for $f_{L,\varepsilon}$ against
$Y_{h_L}$.  By (MA-PI.145)--(MA-PI.147), the pairing
converges to the nonzero pairing obtained after
(MA-PI.143).  A sufficiently small finite positive
$\varepsilon$ supplies $f_Y=f_{L,\varepsilon}$ and
proves (MA-PI.124).  All asserted inputs and outputs
are actual source vectors; no distributional or
nonflat limiting vector replaces this final witness.
\end{proof}

\paragraph{The exact printed Hilbert realization.}
Let $m=\mu_1(C_K^1)>0$ be the actual compact Haar mass.
With the positive complex norm splitting
$g=\nu(e^x)k$, the source definition gives
\[
 \begin{split}
 (T_A\varphi)(\nu(e^x)k)
     &:=|g|^A\Phi_K(\varphi\otimes\mathbf1_{\mathcal O^{S}})(g)
       =\psi(k)^{-1}e^{bx}\mathcal T_S(\varphi)(x),\\
 (U_bF)(\nu(e^x)k)
     &=m^{-1/2}\psi(k)^{-1}e^{bx}F(x),
                    \qquad b=A-\tfrac12 .
 \end{split}
 \tag{MA-PI.148}
\]
These formulas apply to the source vectors and cutoff
outputs used above, which have the same compact type.
The printed polynomial norm makes $U_b$ a unitary
from $H^{\rm one}_{b,s}$ onto that compact-character
component.  Its inverse is compact-fiber integration
against $\psi$, multiplied by $m^{-1/2}e^{-bx}$,
as in (MA-PI.35).  Thus the actual printed orthogonal
projector is conjugate to the one-sided projector
used above, by equality of the full source ranges
under this actual isometry.  Only direct membership
of the test vector in that range was used in the proof.
Multiplying (MA-PI.124) by this unitary gives failure
of the printed canonical intertwining for every
$0<A<1$, in particular the original twist $A=1/2$.
The flat-source error exponent in that realization
is $1+2b=2A$, exactly as in (MA-PI.145).

\paragraph{Sources and Grothendieck placement.}
The original flat source is MA-GA.2 with its inherited
LF topology.  The invariant PW domain is Masullo
Definition~7.3.2.3, (7.3.2.12--17), pp.~354--355,
with $S'=S$ and no extra unit-shell transition set.
Its printed global orthogonal projector is Definition
11.2.1.3, (11.2.1.7--8), pp.~488--489.
The uncentered theta, printed polynomial norm and
actual twist are the source definitions checked in
(MA-PI.32)--(MA-PI.36); their general-number-field
form is unchanged here.  The physical cutoff is
the first Step~II bullet on p.~516 and composition
on p.~517 in Theorem~12.4.1.1.  Our $S$ contains
the field-ramified place $3$, as its p.~514
hypothesis requires.  The complex local gamma
normalization was freshly checked in Tate's actual
Theorem~2.4.1 and complex table, thesis pp.~(2.11)
and~(2.16), PDF pp.~18 and~23.  The new filter,
reflection scale, Gram temperateness and all error
estimates are proved above.

The arithmetic object remains the rank-one finite
Artin character of $\operatorname{Gal}(E/K)=C_3$,
with the established nontrivial two-dimensional
$S_3$ induction to $\mathbb Q$.  At $\mathfrak q$
its inertia invariants vanish; at $\mathfrak p$
it is unramified and has primitive Euler factor
$(1-3^{-s})^{-1}$.  The factor $E_3$ is imposed by
the explicit local test vector and is not a new
arithmetic conductor or a change of motive.
The finite-monodromy convention is SGA~1,
Expos\'e~V, \S7, and Deligne, \emph{Valeurs de fonctions L},
Definition~6.1 and \S\S6.2--6.3, with the inherited
MA-CH/MA-CQ audit scope.  The cote~1 comparison
programme, BR-TR-01/02, scans~079--080 and~111--112,
still requires actual comparison maps before trace
cyclicity; those archive locators are inherited.
This result refutes the specified canonical physical
cutoff/orthogonal projection map in a genuinely
non-induced sector.  It does not assert that the source
literally prints that proposed intertwining equation.
It does not refute a different
Gram-corrected transport, alter the valid nuclear
returns, or remove the arithmetic FT002 condition.
